% Template:     Reporte LaTeX
% Documento:    diffThreshold sweep comparison: pre-fix vs post-fix
% Codificación: UTF-8

\documentclass[
	english,
	letterpaper, oneside
]{article}

\def\documenttitle {Pre-Fix vs.\ Post-Fix \texttt{diffThreshold} Sensitivity}
\def\documentsubtitle {Sanity-Check Repeat of the diffT Sweep on the Original (Buggy) Binary}
\def\documentsubject {Adaptive Subdivision, Load Balancing}
\def\documentdate {\today}

\def\documentauthor {Matias Saavedra}
\def\coursename {Multicore and GPU Programming}
\def\coursecode {CISC 701}

\def\universityname {Universidad de Chile}
\def\universityfaculty {Facultad de Ciencias Físicas y Matemáticas}
\def\universitydepartment {Departamento de Ciencias de la Computación}
\def\universitydepartmentimage {departamentos/dcc}
\def\universitylocation {Santiago de Chile}

../profiling_report/template.tex

\begin{document}

\templatePagecfg
\templateFinalcfg

\inserttitle

% -----------------------------------------------------------------------
\begin{abstract}
This report repeats the \texttt{diffThreshold} sensitivity sweep of
\texttt{difft\_report.tex} on the original (buggy) binary
(\texttt{examine\_rewrite} branch) as a sanity check, and compares the
two runs side by side. The methodology is identical in every other
respect: five threshold values
$\{0.05, 0.1, 0.2, 0.3, 0.5\}$ across three configurations
(\texttt{CPU12}, \texttt{GPU}, \texttt{GPU+11CPU}), three reps each,
plus one \texttt{quiet=0} rep per \texttt{diffT} of \texttt{GPU+11CPU}
for max-region data. The headline finding is that \emph{the optimal
\texttt{diffT} is the same on both branches} --- the hybrid sweet
spot sits at $\texttt{diffT}=0.1$ with means within 0.1~s of each
other ($52.85$~s pre-fix, $52.80$~s post-fix). However, the bug
fix's wall-time delta is strongly \emph{config-dependent}:
\textbf{CPU-only runs $\sim$1\% faster post-fix at every \texttt{diffT}},
while \textbf{GPU-only runs $\sim$1\% slower post-fix at every \texttt{diffT}},
and \textbf{the hybrid is essentially neutral} (within $\pm 0.5\%$)
across the tighter thresholds, becoming $+1.8\%$ slower only at the
legacy default of $0.5$. The bug-fix overhead is real but is exactly
cancelled by improved load balance on configurations where there are
multiple CPU workers to absorb the extra leaves; for the GPU-only
case (one worker, no load-balance benefit available) the overhead
shows up cleanly. The 960$\times$540, $\sim$15~s worst-case CPU region
remains immovable on both branches at every threshold, confirming it
is a heuristic problem and neither a tracking-bug artefact nor a
threshold-tuning issue.
\end{abstract}

% -----------------------------------------------------------------------
\section{Setup}

The binary measured in this report is the one on the
\texttt{examine\_rewrite} branch (the original code, with the
\texttt{else~if} min/max reduction described in
\texttt{bug\_report.tex}); the binary measured in
\texttt{difft\_report.tex} was the
\texttt{examine\_minmax\_bugfix} branch. Everything else is identical:

\begin{itemize}
    \item \textbf{Hardware and workload}: i7-9750H + GTX~1660~Ti~Max-Q,
    100 frames at $1920\times1080$, default
    \texttt{pixelSizeThresh}\,=\,$32{,}768$.
    \item \textbf{Threshold sweep}: $\{0.05, 0.1, 0.2, 0.3, 0.5\}$.
    The $0.5$ baseline for each branch is reused from the corresponding
    Fig~11.15 sweep --- \texttt{sweep\_results/} for pre-fix and
    \texttt{sweep\_results\_postfix/} for post-fix.
    \item \textbf{Configurations and reps}: \texttt{CPU12}, \texttt{GPU},
    \texttt{GPU+11CPU}; 3 reps each.
    \item \textbf{Max-region pass}: 5 single-rep \texttt{quiet=0} runs
    per branch (one per \texttt{diffT}) on \texttt{GPU+11CPU} to
    capture per-region times.
\end{itemize}

% -----------------------------------------------------------------------
\section{Side-By-Side Results}

\subsection{Wall-Clock Times}

\begin{table}[H]
    \centering
    \caption{Mean wall-clock time (s) per config and \texttt{diffT},
             3 reps each. $\Delta$ is post-fix minus pre-fix. Negative
             $\Delta$ means the bug fix \emph{helped} this
             configuration; positive means it \emph{hurt}.}
    \label{tab:wall}
    \begin{tabular}{lrrrrr}
        \hline
        \textbf{Config / diffT} & \textbf{0.05} & \textbf{0.10} & \textbf{0.20} & \textbf{0.30} & \textbf{0.50} \\
        \hline
        CPU12 pre-fix   & 163.12 & 163.90 & 163.88 & 163.58 & 163.47 \\
        CPU12 post-fix  & 161.47 & 161.96 & 162.59 & 162.24 & 162.76 \\
        $\Delta$ (\%)   & $-1.0$\% & $-1.2$\% & $-0.8$\% & $-0.8$\% & $-0.4$\% \\
        \hline
        GPU pre-fix     &  77.27 &  77.18 &  77.14 &  77.14 &  76.54 \\
        GPU post-fix    &  77.60 &  77.81 &  77.82 &  77.83 &  78.03 \\
        $\Delta$ (\%)   & $+0.4$\% & $+0.8$\% & $+0.9$\% & $+0.9$\% & $+1.9$\% \\
        \hline
        GPU+11CPU pre-fix  &  53.24 &  52.85 &  53.19 &  53.30 &  53.62 \\
        GPU+11CPU post-fix &  53.45 &  52.80 &  52.90 &  53.21 &  54.61 \\
        $\Delta$ (\%)      & $+0.4$\% & $-0.1$\% & $-0.5$\% & $-0.2$\% & $+1.8$\% \\
        \hline
    \end{tabular}
\end{table}

\begin{table}[H]
    \centering
    \caption{Population standard deviations (s) per config and
             \texttt{diffT}, 3 reps each.}
    \label{tab:std}
    \begin{tabular}{lrrrrr}
        \hline
        \textbf{Config / diffT} & \textbf{0.05} & \textbf{0.10} & \textbf{0.20} & \textbf{0.30} & \textbf{0.50} \\
        \hline
        CPU12 pre/post      & 1.85 / 0.47 & 0.45 / 0.71 & 0.53 / 0.32 & 0.25 / 0.94 & 0.32 / 2.44 \\
        GPU pre/post        & 0.09 / 0.03 & 0.09 / 0.08 & 0.09 / 0.08 & 0.09 / 0.09 & 0.14 / 0.21 \\
        GPU+11CPU pre/post  & 0.51 / 1.26 & 0.79 / 0.50 & 0.49 / 0.25 & 0.47 / 0.75 & 0.10 / 0.39 \\
        \hline
    \end{tabular}
\end{table}

\insertimage{../difft_sweep_prefix/compare.png}{width=\linewidth}
            {Pre-fix vs.\ post-fix mean wall time per \texttt{diffT} on
             a log-x axis, one panel per configuration; error bars
             show $\pm 1\sigma$. CPU-only is uniformly faster
             post-fix; GPU-only is uniformly slower post-fix; the
             hybrid curve is nearly coincident except at the legacy
             default \texttt{diffT}\,=\,$0.5$. \label{fig:compare}}

\subsection{Leaf Counts Generated}

The bug suppresses splits that the corrected reduction permits, so
the post-fix run always generates more leaves. The gap is largest at
low \texttt{diffT} (where the threshold is tight enough that many
borderline-spread regions need to be reclassified) and narrows at the
legacy default.

\begin{table}[H]
    \centering
    \caption{Total leaves per run (workload-independent at fixed
             \texttt{diffT}: equal across all three configurations on
             the same branch).}
    \label{tab:leaves}
    \begin{tabular}{lrrrrr}
        \hline
        \textbf{diffT} & \textbf{0.05} & \textbf{0.10} & \textbf{0.20} & \textbf{0.30} & \textbf{0.50} \\
        \hline
        Pre-fix leaves   & 4{,}900 & 4{,}900 & 4{,}888 & 4{,}816 & 4{,}603 \\
        Post-fix leaves  & 5{,}323 & 5{,}323 & 5{,}320 & 5{,}290 & 5{,}152 \\
        $\Delta$ leaves  & $+423$ & $+423$ & $+432$ & $+474$ & $+549$ \\
        $\Delta$ \%      & $+8.6$\% & $+8.6$\% & $+8.8$\% & $+9.8$\% & $+11.9$\% \\
        \hline
    \end{tabular}
\end{table}

\subsection{Worst-Case CPU Region (single rep, \texttt{quiet=0})}

\begin{table}[H]
    \centering
    \caption{Per-branch worst-case CPU region from a single
             \texttt{quiet=0} rep of \texttt{GPU+11CPU} at each
             \texttt{diffT}. The size column gives pixel width$\times$height.}
    \label{tab:worst}
    \begin{tabular}{rrlrl}
        \hline
        \textbf{diffT} & \textbf{Pre-fix maxCPU (s)} & \textbf{Pre size} & \textbf{Post-fix maxCPU (s)} & \textbf{Post size} \\
        \hline
        0.05 & 13.32 & $960\times540$ & 14.00 & $960\times540$ \\
        0.10 & 15.31 & $960\times540$ & 15.47 & $960\times540$ \\
        0.20 & 15.13 & $960\times540$ & 15.35 & $960\times540$ \\
        0.30 & 15.35 & $960\times540$ & 15.41 & $960\times540$ \\
        0.50 & 15.62 & $960\times540$ & 15.29 & $960\times540$ \\
        \hline
    \end{tabular}
\end{table}

% -----------------------------------------------------------------------
\section{Analysis}

\subsection{The Optimum \texttt{diffT} Is Branch-Independent}

The single most reassuring result of this sanity check is that the
shape of each configuration's wall-time-vs-\texttt{diffT} curve is
preserved across the fix:

\begin{itemize}
    \item For the hybrid \texttt{GPU+11CPU}, the best value of
    \texttt{diffT} is $0.1$ on both branches, with means
    $52.85$~s pre-fix and $52.80$~s post-fix --- a delta of $50$~ms,
    below noise.
    \item For \texttt{GPU} alone, the best value of \texttt{diffT} is
    $0.5$ on the pre-fix branch ($76.54$~s) but $0.05$ on the post-fix
    branch ($77.60$~s). However, this is a slight artefact of how
    differently the two endpoints behave: see the GPU discussion below.
    \item For \texttt{CPU12}, every threshold lands in a 1.4-second
    band ($161.5$--$162.8$~s post-fix, $163.1$--$163.9$~s pre-fix).
    The CPU-only configuration is essentially insensitive to
    \texttt{diffT} on either branch.
\end{itemize}

\noindent
This is the result a reviewer would want to see: the fix does not
change \emph{which} \texttt{diffT} to pick, only how each \texttt{diffT}
performs. The recommendation from \texttt{difft\_report.tex} to lower
the default from $0.5$ to $0.1$ is robust to the fix.

\subsection{The Bug Fix Is Strictly Beneficial on CPU-Only and Strictly Costly on GPU-Only}

The sign of the post-fix vs.\ pre-fix delta is uniform within each
configuration:

\begin{itemize}
    \item \textbf{CPU12 is always faster post-fix} ($-0.4\%$ to
    $-1.2\%$). The mechanism is direct: with 12 CPU threads, the +423
    extra leaves at \texttt{diffT}\,$\le$\,$0.1$ are absorbed in
    parallel (35 extra leaves per thread on average), and the resulting
    finer-grained work makes load balance across threads tighter. The
    overhead of generating the extra leaves is more than compensated.
    \item \textbf{GPU is always slower post-fix} ($+0.4\%$ to
    $+1.9\%$). With only one GPU thread, every extra leaf is one extra
    kernel launch + memcpy + event-sync round-trip on the same hardware
    serial. There is no parallelism to absorb the overhead; the cost
    shows up directly. The biggest delta is at \texttt{diffT}\,=\,$0.5$
    where the fix adds $+549$ leaves (5{,}152 vs.\ 4{,}603);
    \texttt{GPU}'s wall time rises by $+1.5$~s ($+1.9\%$), almost
    exactly accounted for by the $549$ extra kernel launches at a few
    ms each.
    \item \textbf{GPU+11CPU is essentially neutral across the curve
    except at \texttt{diffT}\,=\,$0.5$.} At every other threshold, the
    post-fix delta lands within $\pm 0.5\%$ of pre-fix. The hybrid
    case has 11 CPU workers \emph{and} a GPU thread, so the GPU's
    overhead penalty and the CPU pool's load-balance benefit partly
    cancel. They cancel less well at \texttt{diffT}\,=\,$0.5$ because
    the absolute extra-leaf count is biggest there ($+549$); they
    cancel essentially perfectly at tighter thresholds.
\end{itemize}

\noindent
This is a sharper version of the finding originally documented in
\texttt{postfix\_report.tex} from the Fig~11.15 sweep: the bug fix's
performance cost depends on which workers are available to absorb the
additional splits. The earlier report observed the slowdown at the
legacy default ($+1.8\%$ on hybrid at \texttt{diffT}\,=\,$0.5$); this
report adds the corresponding picture across the entire threshold
range, and reveals that the slowdown is not an intrinsic property of
the fix but a property of the threshold value at which the fix was
measured.

\subsection{The 960$\times$540 Outlier Is Immune on Both Branches}

The pixel dimensions of the worst single CPU region in every column of
Table~\ref{tab:worst} are $960\times540$. Both branches, every
\texttt{diffT}. The compute time of this region ranges 13.3--15.6~s
across the ten cells of the table, mostly within noise of each other.
The bug fix did not unlock this outlier (it could not have ---
\texttt{bug\_report.tex} explained why the bug discards corner-0's
maximum only, while this region has all four corners at
\texttt{MAXITER} so its corner spread is identically zero, beyond
either branch's classifier's reach). The diffT sweep did not unlock
this outlier (a tighter threshold still cannot reject zero spread).
Two independent levers, both with no effect on the binding
constraint.

This is the cleanest possible evidence that the worst-case region is
a heuristic problem and not a tracking-bug or tuning artefact.

\subsection{Variance: A Subtler Story}

Standard deviations tell an unexpected story
(Table~\ref{tab:std}):

\begin{itemize}
    \item \textbf{\texttt{CPU12} at \texttt{diffT}\,=\,$0.5$}: pre-fix
    $\sigma = 0.32$~s, post-fix $\sigma = 2.44$~s. The post-fix
    variance is $8 \times$ higher at the legacy default. The
    \texttt{difft\_report.tex} attribution was thread-scheduling
    jitter on the worst-case region; the pre-fix data shows that this
    jitter is much smaller on the buggy branch.
    \item \textbf{\texttt{GPU+11CPU} at \texttt{diffT}\,=\,$0.5$}:
    pre-fix $\sigma = 0.10$~s, post-fix $\sigma = 0.39$~s. Again
    $\sim 4 \times$ noisier post-fix at the legacy default.
    \item \textbf{\texttt{GPU+11CPU} at \texttt{diffT}\,=\,$0.05$}:
    pre-fix $\sigma = 0.51$~s, post-fix $\sigma = 1.26$~s.
    $2.5 \times$ noisier post-fix at the most aggressive value.
    \item \textbf{\texttt{CPU12} at \texttt{diffT}\,=\,$0.05$}: pre-fix
    $\sigma = 1.85$~s, post-fix $\sigma = 0.47$~s. Here the pattern
    \emph{inverts}: the post-fix is more deterministic.
\end{itemize}

\noindent
A consistent interpretation: the buggy classifier sometimes makes the
same misclassifications across reps (since the bug depends on the
\emph{order} in which corner values are visited, which is
deterministic from the geometric subdivision tree); these systematic
misclassifications stabilise some configurations into more repeatable
schedules. Fixing the bug exposes the workload's underlying
non-deterministic scheduling sensitivities, which then become visible
as run-to-run variance. The exception is \texttt{CPU12} at
\texttt{diffT}\,=\,$0.05$: at the most aggressive threshold on the
buggy branch there are not enough leaves (4{,}900) to amortise the
worst-case region across 12 threads, and the worst-region scheduling
becomes unstable; the fix's $+423$ leaves give the queue enough
material to dilute the variance.

\subsection{Why the Pre-Fix Curve Has a Different Best for GPU}

On the pre-fix branch, the \texttt{GPU} configuration is best at
\texttt{diffT}\,=\,$0.5$ ($76.54$~s), with the four-tightest
thresholds clustered between $77.1$ and $77.3$~s. This is the opposite
ordering from what one might naively expect: the buggy code's
worst-classification \emph{helps} the single GPU thread because
\emph{not} splitting reduces the per-launch overhead the GPU thread
pays. The fix recovers $+549$ leaves at \texttt{diffT}\,=\,$0.5$ and
those translate directly into a $+1.5$~s wall-time penalty for the
GPU-only configuration. The post-fix branch's monotonic improvement
from $0.5$ down to $0.05$ shows that, with the bug corrected,
\emph{less aggressive thresholds} are now the worse choice for
GPU-only.

This is a useful confirmation that the fix's per-launch overhead is
real and quantifiable.

% -----------------------------------------------------------------------
\section{Summary of Recommendations}

The pre-fix sanity check confirms every actionable conclusion from
\texttt{difft\_report.tex}:

\begin{enumerate}
    \item \textbf{Use \texttt{diffT}\,=\,$0.1$ as the default.} Both
    branches agree this is the hybrid sweet spot. The fix-induced
    overhead at \texttt{diffT}\,=\,$0.5$ disappears at this
    tighter value.
    \item \textbf{Avoid \texttt{diffT}\,=\,$0.05$.} Post-fix variance
    is unacceptably high there; pre-fix variance is bad too. The
    pixel-size floor saturates the splitting before this threshold
    can contribute additional useful work.
    \item \textbf{The corner heuristic itself is the bottleneck.}
    Two orthogonal experiments (the diffT sweep on each branch)
    independently fail to subdivide the $960\times540$, $\sim$15~s
    worst-case region. The only paths forward are edge-midpoint
    sampling, cardioid/bulb certificates, or async streams; none of
    these can be replaced by parameter tuning.
\end{enumerate}

% -----------------------------------------------------------------------
\section{Conclusion}

Repeating the \texttt{diffThreshold} sensitivity sweep on the
original buggy binary produces a result that is reassuring along the
dimensions that matter:

\begin{itemize}
    \item The optimal \texttt{diffT} is the same on both branches.
    \item The shape of each wall-time curve is preserved.
    \item The bug-fix's wall-time delta is exactly the per-configuration
    cost-or-benefit one would predict from the load-balance argument
    in \texttt{postfix\_report.tex}: cost on GPU (no parallelism to
    absorb overhead), benefit on CPU (many threads to absorb extra
    leaves), and approximate neutrality on hybrid (the two effects
    cancel).
    \item The worst-case region is invariant to both \texttt{diffT}
    and the bug fix.
\end{itemize}

\noindent
The sanity check uncovered no contradiction with the earlier
reports, but it did sharpen one quantitative claim: the
$+1.8\%$ hybrid slowdown observed in \texttt{postfix\_report.tex} at
the legacy default \texttt{diffT}\,=\,$0.5$ is the worst case of the
post-fix performance penalty across the entire \texttt{diffT} range
that has been tested. At the recommended new default of $0.1$, the
penalty is zero within noise. The bug fix is therefore a correctness
improvement with a tuning-dependent performance cost, and at the
right \texttt{diffT} that cost vanishes.

\end{document}
